Compliance auditing against ISO/IEC 27001:2022 requires mapping organizational policy documents to specific Annex A controls, a process that is time-consuming and prone to inconsistency when performed manually. This study develops and evaluates information retrieval strategies within a Retrieval-Augmented Generation (RAG) framework to support automated Statement of Applicability (SoA) control mapping for ISO/IEC 27001:2022. Three design axes are systematically evaluated: (1) the weight composition between dense and sparse retrieval in hybrid retrieval, (2) the candidate pool size for cross-encoder reranking, and (3) the score fusion weight between hybrid retrieval and cross-encoder scores.

Experiments are conducted using 77 queries derived from four policy documents of the Directorate of Information Technology (DTI), Universitas Gadjah Mada, against a knowledge base of 93 Annex A controls (A.5 through A.8) from ISO/IEC 27001:2022. The knowledge base is structured as JSON entities containing descriptive attributes, audit indicators, keywords, related assets, and security principles in both English and Indonesian. The sparse retrieval component employs BM25, the dense retrieval component uses the paraphrase-multilingual-MiniLM-L12-v2 embedding model, and the reranking component employs the BAAI/bge-reranker-v2-m3 cross-encoder. Evaluation metrics include Hit@3, Recall@3, Mean Average Precision (MAP@3), and NDCG@3, supplemented by non-parametric Friedman statistical testing.

Experimental results demonstrate that the balanced hybrid retrieval configuration with dense=0.50 and sparse=0.50 yields the best performance on three of four metrics (Recall@3=0.537; MAP@3=0.467; NDCG@3=0.575). A candidate pool size of 5 achieves the highest reranking quality with a monotonic performance decline as pool size increases. The hybrid-dominant score fusion with Hybrid:CE=0.75:0.25 provides the best overall performance (Recall@3=0.541; NDCG@3=0.579). The Friedman test confirms that performance differences across configurations are statistically significant ($p < 0.05$) in all experimental scenarios. A per-query analysis reveals that hybrid retrieval and cross-encoder reranking exhibit complementary strengths: hybrid retrieval excels on queries containing explicit technical terminology (36 out of 77 queries), while cross-encoders benefit queries requiring deeper contextual understanding (15 out of 77 queries).

This study concludes that the optimal configuration of balanced hybrid retrieval, a small candidate pool, and hybrid-dominant score fusion achieves the best trade-off between retrieval quality and computational efficiency for ISO/IEC 27001:2022 SoA control mapping.

\noindent\textbf{Keywords} : ISO/IEC 27001, Statement of Applicability, hybrid retrieval, cross-encoder reranking, retrieval-augmented generation, information security audit
